\documentclass[11pt,a4paper]{article}

\usepackage[margin=1.05in]{geometry}
\usepackage{amsmath,amssymb,amsthm}
\usepackage{graphicx}
\usepackage{mathtools}
\usepackage{booktabs}
\usepackage{xcolor}
\usepackage[colorlinks=true,linkcolor=black,citecolor=black,urlcolor=blue]{hyperref}
\usepackage{enumitem}
\usepackage[T1]{fontenc}
\usepackage[protrusion=true,expansion=false]{microtype}
\setlength{\emergencystretch}{2.5em}

\theoremstyle{plain}
\newtheorem{theorem}{Theorem}[section]
\newtheorem{proposition}[theorem]{Proposition}
\newtheorem{lemma}[theorem]{Lemma}
\newtheorem{conjecture}[theorem]{Conjecture}
\theoremstyle{definition}
\newtheorem{definition}[theorem]{Definition}
\newtheorem{example}[theorem]{Example}
\newtheorem{decision}[theorem]{Decision}
\newtheorem{question}[theorem]{Open question}
\theoremstyle{remark}
\newtheorem{remark}[theorem]{Remark}

%%%%%%%%%%%%%%%%%%%%%%%%%%%%%%%%%%%%%%%%%%%%%%%%%%%%%%%%%%%%%%%%%%%%%%%%%%%%%%%
% NOTATION RULES OBSERVED THROUGHOUT:
%   * the rho calculus is never written with the Greek letter rho
%   * quoting is @P, dropping is *x; corner quotes are pre-2016 notation
%%%%%%%%%%%%%%%%%%%%%%%%%%%%%%%%%%%%%%%%%%%%%%%%%%%%%%%%%%%%%%%%%%%%%%%%%%%%%%%
\newcommand{\rhoc}{\textsf{rho}}
\newcommand{\sem}[1]{[\![#1]\!]}
\newcommand{\ket}[1]{\lvert #1 \rangle}
\newcommand{\cpl}[1]{\langle\!\langle #1 \rangle\!\rangle}
\newcommand{\Zt}{\mathbb{Z}/2}
\newcommand{\Ze}{\mathbb{Z}/8}
\newcommand{\Sc}{\mathbb{Z}[\tfrac{1}{\sqrt2},\omega]}
\newcommand{\out}[2]{#1!#2}
\newcommand{\rcv}[3]{\mathsf{for}(#1 \leftarrow #2)#3}
\newcommand{\drop}[1]{{*}#1}
\newcommand{\bdrop}[2]{{*}(#1,\,#2)}
\newcommand{\Sig}[2]{\Sigma #1.\,#2}
\newcommand{\red}{\rightarrow}
\newcommand{\scong}{\equiv}
\newcommand{\fv}{\mathrm{fv}}

\title{Unitary in the Equations, Measurement in the Reductions\\[0.4em]
\large A paired-message \rhoc{} calculus with a value binder,\\
path-sum equations, and traced execution}
\author{L.G. Meredith\\ F1R3FLY.io}
\date{Working note, 1 October 2026}

\begin{document}
\maketitle

\begin{abstract}
This note records a design session exploring a single intuition: in a
graph-structured lambda theory the \emph{unitary} behaviour lives in the
(possibly augmented) equations, while \emph{measurement} lives in the
reductions. The \rhoc{} calculus with graded where-clauses is the vehicle,
because every weighted GSLT translates into it. Working the intuition out
forces a small number of changes to \rhoc{}: messages and names become
\emph{pairs} of a process and a phase label; dropping splits into a binary,
phase-preserving \emph{equation} and a unary, semantically-substituted
\emph{execution} whose erasure is \emph{traced}; superposition enters through
a binder $\Sigma$ over $\Zt$-valued variables; and couplings between values are
inert atoms governed by path-sum equations, which are the algebraic face of the
graph-like ZX-calculus. Reductions then come in exactly two measuring flavours
(resolving a value, or running received code) and one coherent flavour
(binding receipts). Persistent receipts are dropped, recursion is recovered by
reflection, and loop guards are classical by theorem. A worked example computes
$HZH\ket{0} = \ket{1}$ by equations alone, and contrasts it with a
which-rail detector that destroys the interference by reductions alone. The
note closes with the decisions taken and the questions left open.
\end{abstract}

\tableofcontents

%=============================================================================
\section{The intuition and where it comes from}
%=============================================================================

The intuition under study is:

\begin{quote}
\emph{The unitary behaviour of a quantum-interpreted GSLT lives in its
(possibly augmented) equations; measurement lives in its reductions.}
\end{quote}

It has roots in earlier work in this line. The first version of the
weighted-GSLT note assigned the Hamiltonian to equations and the jump
operators to rewrites, and its second version distilled the slogan that a
weighted GSLT is quantum exactly to the extent that its presentation withholds
equations from the quotient. What is new here is taking the intuition as the
\emph{organising principle} for \rhoc{} with graded where-clauses. That
reverses the direction of the graded-where programme
\cite{GWv6}, in which reductions carried coherent amplitude: here reductions
are where coherence ends.

Three observations make the reversal natural in \rhoc{}.
\begin{itemize}[leftmargin=1.5em]
\item \textbf{COMM quotes.} Firing substitutes $@Q$ into a continuation; the
  payload becomes a name, and names are compared by equality. A communication
  therefore writes a record, which is precisely Barandes's criterion for a
  division event \cite{Barandes}.
\item \textbf{Cost.} Equations are free and rewrites are costed. Cost lands
  where information is erased, which is Landauer's principle read
  operationally.
\item \textbf{Hermiticity.} An equation is symmetric; a directed rewrite is
  not. A weighted symmetric relation can be Hermitian; a directed one cannot.
\end{itemize}

\paragraph{Reading order.} Section~\ref{sec:path} retraces the design path,
because several of the choices only make sense against the alternatives that
were rejected. Section~\ref{sec:calculus} states the calculus.
Section~\ref{sec:props} records its basic properties.
Section~\ref{sec:example} works an example. Sections~\ref{sec:context}
and~\ref{sec:open} place the work and list what remains.

%=============================================================================
\section{The design path}\label{sec:path}
%=============================================================================

\subsection{Embedding ZX: the graph-like fragment}

The starting point was to embed the ZX-calculus \cite{CoeckeDuncan} into
\rhoc{}, and the cleanest target is the \emph{graph-like} fragment
\cite{DKPW}: only green spiders, joined by Hadamard edges, into which every
diagram rewrites and to which PyZX normalises. Reading a quantum name as a
$Z$-spider, scope as the boundary, and AC of parallel composition plus scope
extrusion as ``only connectivity matters'' gives spider fusion and identity
removal for free.

Two lessons came from this first pass and survive into the final design.
\begin{enumerate}[leftmargin=1.5em]
\item \textbf{Once the congruence contains a complete ZX theory, syntactic
  matching on quantum structure is not well defined.} The $\scong$-class of a
  quantum subterm is its linear map, so the only $\scong$-invariant
  observations are semantic. Deciding $\scong$ is simulation.
\item \textbf{Unbinding is post-selection.} Making a spider internal plugs
  $\langle 0\rvert + \langle 1\rvert = \sqrt2\langle +\rvert$ into it.
  Genuine discard needs the ground of ZX$_{\perp}$ \cite{CJPV}, and that is
  where reductions belong.
\end{enumerate}

\subsection{Quoted pairs, and closure under reflection}

The proposal that channels be quoted pairs $@(P,\phi)$ moves the phase into the
name. Making only channels into pairs breaks the invariant that keeps \rhoc{}
closed under reflection: the payload sort and the quoted sort must coincide.
Making \emph{messages} pairs as well, $x!(Q,\phi)$, restores it. A
communication substitutes $@(Q,\phi)$, which is again a name.

\subsection{Binary drop}

The remaining seam was the drop in process position, which would yield a pair
where a process is expected. The resolution is a \emph{binary} drop
\[
  \bdrop{@(P,\phi)}{u} \;:=\; u!(P,\phi),
\]
which never produces a process. It is invertible (the message determines both
arguments), it preserves the phase, and it is an equation. On its own it gives
an exactly closed calculus, but nothing executes quoted code any more, so
recursion is lost; with linear receipts every term strongly normalises.

\subsection{Semantic substitution and traced erasure}

Recursion is recovered as in the original \rhoc{} paper \cite{MR05}: drop in
process position is resolved by \emph{semantic substitution}, and on
substitution the process begins running and its phase is erased.

Erasure read literally, $(P,\phi)\mapsto P$, is many-to-one; amplitudes of
branches that differ only in label add, which is not an isometry and
reintroduces the norm leak (and with it PostBQP). Erasure is therefore read as
a \emph{trace}: the information goes to an environment nobody reads, and the
step is a completely positive instrument. Substitution is \emph{eager}: the
trace happens at the communication, even if the drop is guarded and never
reached.

What the trace records was refined twice. Recording only the label lets
superposed control through (two routed payloads with equal labels but
different code). Recording the \emph{whole message} closes the gap, and is the
choice most consistent with context-labelled bisimulation \cite{LeiferMilner},
where what an observer learns from a transition is its label and labels carry
whole contexts. Messages are compared up to structural congruence.

\subsection{The coupling: bicharacter, and by equation}

At a communication, a sending label $a$ and a receiving label $b$ could combine
as a difference $e^{i(a-b)}$ or as a bicharacter $\chi(a,b)$. The difference is
a $U(1)$ connection with gauge freedom per channel, and it has a sharp
limitation: every maximal matching of a complete join uses each occurrence
once, so all matchings carry the same total phase and the site matrix has rank
one. The bicharacter is what ZX is built on. Hadamard has entries
$(-1)^{ab}/\sqrt2$, a Hadamard edge in graph-like ZX is a $CZ$ with amplitude
$(-1)^{ab}$, and a $Z$-spider of phase $k\pi/4$ contributes $\omega^{ka}$
with $\omega = e^{i\pi/4}$, which is a bicharacter
$\Ze\times\Zt\to U(1)$. In ZX the coupling is static, so it is evaluated by
\emph{equation}, not at the communication.

\subsection{A correction: labels are data, couplings are atoms}

A static coupling needs values to sum over, hence a binder. The first attempt
put value-dependent terms into the phase slot of a pair. That fails for two
reasons.
\begin{itemize}[leftmargin=1.5em]
\item A label that depends on a summed variable is \emph{data in the
  configuration}, so the branches it separates are distinct basis
  configurations. It is a which-path record, and it destroys the very
  interference the coupling was meant to produce.
\item The trace reads a message's label, and the path-sum equations rewrite
  labels, so the reduction would not be invariant under $\scong$.
\end{itemize}
The resolution separates the two roles. The phase slot of a pair holds a
\emph{constant} $k\in\Ze$, which is data: part of the name, and traced on
execution. Couplings between values live in a separate inert atom
$\cpl{\phi}$ that is never received, run or traced, so the equations may
rearrange it freely.

%=============================================================================
\section{The calculus}\label{sec:calculus}
%=============================================================================

\subsection{Syntax}

\begin{definition}[Syntax]\label{def:syntax}
\[
\begin{array}{rcll}
e & ::= & v \mid 0 \mid 1 \mid e\oplus e & \text{values, in } \Zt\\[2pt]
\phi & ::= & k\cdot m \mid \phi+\phi & k\in\Ze,\ m \text{ a monomial in values}\\[2pt]
s & \in & \Sc & \text{scalars}, \ \omega=e^{i\pi/4}\\[2pt]
M & ::= & (P,k) & \text{messages}, \ k\in\Ze \text{ constant}\\[2pt]
x & ::= & @M \mid y \mid x[e] & \text{names}\\[2pt]
P & ::= & 0 \mid x!M \mid \rcv{y}{x}{P} \mid \rcv{y}{x[u]}{P} & \\
  & \mid & \drop{y} \mid \bdrop{x}{u} \mid P \mid P
        \mid \Sig{v}{P} \mid \cpl{\phi} \mid s\cdot P & \text{processes}
\end{array}
\]
The receipt $\rcv{y}{x[u]}{P}$ binds both $y$ and the index variable $u$.
$\Sig{v}{P}$ binds the value variable $v$ in $P$, including occurrences in
indices and couplings. For Clifford, monomials of degree at most two suffice;
$(-1)^{ab}$ is written $\cpl{4\cdot ab}$.
\end{definition}

\begin{definition}[Name identity]\label{def:identity}
Indexed names $x[0]$ and $x[1]$ are distinct. A name $@(P,k)$ is identified by
its whole pair; the label $k$ is part of the name, as data. A name $x[e]$ with
$e$ depending on $\Sigma$-bound variables denotes different channels in
different branches.
\end{definition}

\begin{remark}[What is not primitive]
There are no persistent receipts, no replication and no summation operator on
processes. There is no restriction operator either; fresh names arise by
quotation. The binder $\Sigma$ is deliberately separate from any freshness
binder: $\Sigma$ gives a spider a value, a fresh name gives a channel an
identity, and conflating them would make every fresh channel a superposition.
\end{remark}

\subsection{Equations: the unitary side}

\begin{definition}[Structural congruence]\label{def:cong}
$\scong$ is the least congruence, closed under $\alpha$-conversion of both
name and value binders, satisfying:
\begin{description}[leftmargin=2.5em,style=nextline]
\item[(E1) Monoid and scalars] $P\mid 0\scong P$, $P\mid Q\scong Q\mid P$,
  $(P\mid Q)\mid R\scong P\mid(Q\mid R)$; $(s\cdot P)\mid Q\scong s\cdot(P\mid
  Q)$, $s\cdot(t\cdot P)\scong (st)\cdot P$, $1\cdot P\scong P$.
\item[(E2) Binary drop] $\bdrop{@M}{u}\scong u!M$.
\item[(E3) Scope] $(\Sig{v}{P})\mid Q\scong \Sig{v}{(P\mid Q)}$ when
  $v\notin\fv(Q)$; $\Sig{v}{\Sig{w}{P}}\scong\Sig{w}{\Sig{v}{P}}$;
  $\Sig{v}{(s\cdot P)}\scong s\cdot\Sig{v}{P}$.
\item[(E4) Couplings] $\cpl{\phi}\mid\cpl{\psi}\scong\cpl{\phi+\psi}$,
  $\cpl{0}\scong 0$, $\cpl{k\cdot 1}\scong \omega^k\cdot 0$, with arithmetic in
  $\Ze$ and $v^2 = v$.
\item[(E5) Elim] $\Sig{v}{P}\scong 2\cdot P$ when $v\notin\fv(P)$.
\item[(E6) HH] $\Sig{w'}{\Sig{w}{(\cpl{4\cdot w\cdot(w'\oplus e)}\mid P)}}
  \scong 2\cdot P\{e/w'\}$ when $w\notin\fv(P,e)$ and $w'\notin\fv(e)$.
\item[(E7) $\omega$-rule] The rule completing the Clifford fragment of the
  path-sum calculus \cite{Amy}; its statement in this syntax is open
  (Question~\ref{q:e7}).
\end{description}
\end{definition}

E6 is the identity $\sum_{w\in\Zt}(-1)^{w f} = 2\,\delta_{f=0}$: a linear
constraint on a summed variable becomes a substitution. In ZX it is the Hopf
law and Hadamard cancellation; in this calculus it is name fusion. E2 is the
only equation that touches communication syntax, and it is a bijection between
term shapes. Every equation is an equality, hence invertible, so the unitary
content sits where the intuition places it.

\subsection{Denotation}

\begin{definition}[Path-sum denotation]\label{def:denot}
Every closed term is $\scong$ to one of the form
$\Sig{\vec v}{(s\cdot\cpl{\phi}\mid C)}$ with $C$ coupling-free and
scalar-free. It denotes
\[
  \sem{\Sig{\vec v}{(s\cdot\cpl{\phi}\mid C)}}
  \;=\; s\sum_{\vec v\in\Zt^n}\omega^{\phi(\vec v)}\,\ket{C(\vec v)},
\]
where the configurations $C(\vec v)$, taken modulo E1--E3, form an orthonormal
basis.
\end{definition}

This is the path-sum representation of Dawson et al.\ \cite{Dawson} and Amy
\cite{Amy}: the algebraic twin of ZX. Note that in this architecture a term
class \emph{is} a vector; there is no time evolution generated by the
equations, and gates are applied by composition.

\subsection{Reductions: the measurement side}

Reductions are closed under parallel composition, scalars, $\Sigma$, and
$\scong$. Write $\mathrm{idx}(c[e]) = e$ and $\mathrm{idx}(x) = 0$ for an
unindexed name.

\begin{definition}[Reduction rules]\label{def:red}
\begin{description}[leftmargin=2.5em,style=nextline]
\item[(R1) Coherent COMM] The rule is
  \[
  c[e]!(P,k)\mid\rcv{y}{c[u]}{Q}\;\red\;\cpl{k\cdot e}\mid Q\{@(P,k)/y,\ e/u\}.
  \]
  The receipt binds the index symbolically, so nothing is resolved. The
  payload's label is \emph{banked} as amplitude against the index: this is the
  bicharacter coupling of a phase with a value. The unindexed case,
  $x!(P,k)\mid\rcv{y}{x}{Q}\red Q\{@(P,k)/y\}$, banks nothing.
\item[(R2) Resolving COMM]
  A receipt on $c[b]$ with $b$ a constant, facing $c[e]!M$ with $e$ dependent
  on $\Sigma$-bound variables, fires only on the assignments where $e = b$.
  The step is a projective instrument with two outcomes: \emph{fired}, which
  restricts the sum to $e=b$ and then proceeds as R1; and \emph{not fired},
  which restricts to $e\neq b$. Outcome probabilities are Born weights of the
  restricted sums. In general: \emph{a communication whose enabledness varies
  across assignments measures the variables it depends on.}
\item[(R3) Traced execution]
  If $\drop{y}$ occurs in process position anywhere in $Q$ (guarded
  occurrences included), the communication additionally records the received
  message $M=(P,k)$, up to $\scong$, to an environment, and
  \[
  (\drop{y})\{@(P,k)/y\} \;=\; P.
  \]
  On binary drops the substitution is syntactic:
  $\bdrop{y}{u}\{@M/y\} = \bdrop{@M}{u}\scong u!M$.
\end{description}
\end{definition}

\begin{remark}[R3 as an instrument]
With Kraus operators indexed by $\scong$-classes of messages,
$K_{[M]}$ projects onto the branches in which $M$ was consumed and then
substitutes; $\sum K^\dagger K = I$, so the step is trace preserving. Its
Stinespring form writes $M$ to a channel nobody holds. Three perspectives
coincide with those requested of the heralds in \cite{GWv6}: globally the state
is pure; for anyone without the record's capability it is dephased; for anyone
who reads the record it is a measurement with an outcome.
\end{remark}

\begin{remark}[When the trace is inert]
The record is inert, and the step coherent, exactly when the same message (up
to $\scong$) is consumed in every branch. In particular running code with a
constant label and no $\Sigma$-bound values in it costs no coherence.
\end{remark}

%=============================================================================
\section{Properties}\label{sec:props}
%=============================================================================

\begin{proposition}[Closure under reflection]
Quotation takes messages to names and the binary drop takes names back to
messages: $\bdrop{@M}{u}\scong u!M$, and every communication substitutes a name
of the form $@M$. The payload sort and the quoted sort coincide.
\end{proposition}

\begin{proposition}[Replication is derivable]
Let $R := \rcv{y}{x}{(\bdrop{y}{x}\mid\drop{y})}$ and
$!P := x!(R\mid P,0)\mid R$. Then $!P\red x!(R\mid P,0)\mid R\mid P$.
The trace at each unfolding records the same message in every branch and is
therefore inert.
\end{proposition}

\begin{proposition}[Classical control]\label{prop:control}
Within the scope of a $\Sigma$ binder, \emph{which communications fire and
which code runs} is the same under every assignment of the bound values,
except at R2 and R3 events.
\end{proposition}

\begin{proof}[Proof sketch]
$\Sigma$-bound values occur only in indices, couplings and inside quoted
payloads. A message's presence cannot depend on a value except through its
index, so unindexed communications are enabled uniformly, and binding receipts
$c[u]$ fire uniformly. A communication whose enabledness varies is R2, which
resolves; this includes matching a quoted name whose identity depends on a
value. The code that runs can differ across branches only if the messages
consumed differ, and R3 records them.
\end{proof}

\begin{remark}[Consequences of Proposition~\ref{prop:control}]
The coherent part of any run is a fixed circuit whose shape is classical. This
is Selinger's \emph{quantum data, classical control} \cite{Selinger}. Loop
guards are classical, so loops are monotone and the Löwner-order obstruction to
recursion under quantum alternation \cite{BP} does not arise. Mach--Zehnder
style routing survives, because it moves data along rails without forking
code, so the positioning against quantum alternation in \cite{GWv6} is
unaffected.
\end{remark}

\begin{remark}[The discipline the trace imposes]
Running code that mentions a $\Sigma$-bound value measures that value, because
a context could tell the variants apart. Quantum data must therefore enter code
through binding receipts, never through the code itself. Gates written as in
Section~\ref{sec:example} obey this: their code mentions only their rails, and
they acquire the index at the communication.
\end{remark}

\begin{remark}[Measurement criterion]
A communication is coherent when its subject does not depend on a bound value
(or is matched by a binding receipt) \emph{and} its continuation does not run
what it receives (or runs the same message in every branch). Otherwise it
measures. This is a weak reading of ``measurement in the reductions'' with a
syntactic criterion; the strong reading is the special case in which every
continuation runs what it receives.
\end{remark}

%=============================================================================
\section{A worked example}\label{sec:example}
%=============================================================================

A qubit is a token $T := 0$ on a rail $q[v]$. A Hadamard gate is a process:
\[
  H(q,q') \;:=\; \rcv{y}{q[u]}{\;\Sig{w}{\big(2^{-1/2}\cdot\cpl{4\cdot u\cdot w}
  \;\mid\; \bdrop{y}{q'[w]}\big)}}.
\]
It binds the input index $u$, sums a fresh output value $w$ with coupling
$(-1)^{uw}$, and forwards the pair unchanged by binary drop. Nothing is run, so
the gate is coherent.

\begin{example}[$HZH\ket{0}=\ket{1}$ by equations alone]\label{ex:hzh}
Start from
\[
  q_0[0]!(T,4)\;\mid\;H(q_0,q_1)\;\mid\;H(q_1,q_2).
\]
The token carries label $4$, so every hop banks $(-1)^{\mathrm{index}}$; it
acts as a $Z$ at each hop.
\begin{enumerate}[leftmargin=1.5em]
\item \emph{Fire $H(q_0,q_1)$.} By R1, $u:=0$ and $\cpl{4\cdot 0}\scong 0$:
  \[ 2^{-1/2}\cdot\Sig{w}{\,q_1[w]!(T,4)}\;\mid\;H(q_1,q_2), \]
  which denotes $\ket{+}$.
\item \emph{Fire $H(q_1,q_2)$.} After extruding $\Sigma w$ over the gate (E3),
  R1 binds $u:=w$ symbolically and banks $\cpl{4\cdot w}$, which is the $Z$:
  \[ 2^{-1/2}\cdot\Sig{w}{\Big(\cpl{4\cdot w}\mid 2^{-1/2}\cdot
     \Sig{w'}{\big(\cpl{4\cdot w\cdot w'}\mid q_2[w']!(T,4)\big)}\Big)}. \]
\item \emph{Equations only.} By E4 the couplings fuse to
  $\cpl{4\cdot w\cdot(w'\oplus 1)}$; by E1 and E3 the scalars float out and the
  binders reorder to $\Sigma w'.\Sigma w$; E6 with $e=1$ gives
  \[ 2^{-1/2}\cdot 2^{-1/2}\cdot 2\cdot q_2[1]!(T,4) \;=\; q_2[1]!(T,4). \]
\end{enumerate}
The $w'=0$ branch is cancelled by interference inside the equations, with no
reduction involved. A receipt on $q_2[1]$ now fires with probability~$1$.
With label $0$ the same steps give $q_2[0]!(T,0)$, which is $HH=I$.
\end{example}

\begin{example}[A which-rail detector destroys the interference]
Replace the second gate's input rail by detectors:
\[
  \rcv{y}{q_1[0]}{\bdrop{y}{d[0]}}\;\mid\;\rcv{y}{q_1[1]}{\bdrop{y}{d[1]}}
  \;\mid\; H(d,q_2).
\]
Each detector is enabled on only one value of $w$, so by R2 each resolves $w$.
The state becomes an equal mixture of $d[0]!(T,4)$ and $d[1]!(T,4)$, and the
gate then yields each output with probability $1/2$. The equations are
unchanged; only the reductions differ, and they alone remove the interference.
\end{example}

%=============================================================================
\section{Context}\label{sec:context}
%=============================================================================

\paragraph{ZX and path sums.} The equations are path-sum rewrite rules
\cite{Dawson,Amy}, whose completeness for Clifford is known, and which are the
algebraic counterpart of graph-like ZX \cite{DKPW} with its completeness
results \cite{JPV,Vilmart}. R3 plays the role of the ground in ZX$_\perp$
\cite{CJPV}.

\paragraph{Measurement-based computation.} Graph states with local phases plus
single-qubit measurements and classical feed-forward are universal
\cite{RB}, and gflow \cite{BKMP} is the condition under which internal spiders
are realised by real measurements and corrections rather than post-selection.
The calculus has this shape: static couplings, measuring reductions, classical
control. This suggests that a strong reading of the intuition costs no
expressiveness, but that has not been checked here.

\paragraph{Quantum control versus quantum data.} The classical-control theorem
places the calculus with Selinger \cite{Selinger} and with CQP \cite{CQP} as
regards control, while the coherent routing of R1 keeps the Mach--Zehnder
reading of quantum alternation \cite{BP,YYF} available for data.

\paragraph{Barandes.} Division events \cite{Barandes} are R2 and R3 steps:
moments at which the branch distinction is written into a name or a record.
Between them the process is indivisible in his sense.

\paragraph{Context-labelled bisimulation.} R3 records exactly what a
context-labelled transition exposes, so decoherence classes and observation
classes coincide by construction (up to the $\scong$-versus-$\sim$ gap noted in
Question~\ref{q:sim}).

\paragraph{The catamorphism notes.} The weighting notes in the repository
dilate the rewrite operator and post-select per step. This note moves the
unitary out of the rewrites altogether; if the weighting functor is required to
respect the equations, augmentation makes it a projective (twisted) functor.

%=============================================================================
\section{Decisions and open questions}\label{sec:open}
%=============================================================================

\subsection{Decisions taken}

\begin{center}
\begin{tabular}{@{}p{0.34\textwidth}p{0.58\textwidth}@{}}
\toprule
\textbf{Question} & \textbf{Decision}\\
\midrule
Payload sort & Messages are pairs $(P,k)$; names are quoted messages.\\
Drop & Binary drop is an equation; unary drop is semantic substitution.\\
Erasure on execution & Traced, not coherent.\\
When the trace happens & Eagerly, at the communication.\\
What the trace records & The whole message, up to $\scong$.\\
How labels couple & By bicharacter, evaluated by equation.\\
Label versus coupling & Labels are constant data; couplings are inert atoms.\\
Superposition & A dedicated binder $\Sigma$ over $\Zt$.\\
Persistent receipts & Dropped; replication is derivable.\\
Loop guards & Classical (now a theorem).\\
\bottomrule
\end{tabular}
\end{center}

\subsection{Open questions}

\begin{question}[Per-hop banking]\label{q:hop}
As written, R1 banks a payload's label at every indexed channel it crosses, so
a labelled payload acts as a phase gate at every hop. This is consistent but
may be more than intended. The alternative is to bank once, for instance when
the label is first placed.
\end{question}

\begin{question}[Born normalisation]
How normalisation interacts with maximal progress, and whether R2's
``not fired'' outcome should be an outcome at all or be absorbed by maximal
progress.
\end{question}

\begin{question}[Race normalisation]
Re-deriving contention-set completion \cite{GWv6} in this calculus, now that
some communications are instruments and the coherent ones only move pairs.
\end{question}

\begin{question}[E7 and the zero scalar]\label{q:e7}
Stating the $\omega$-rule in this syntax, and deciding how to treat $0\cdot P$,
which collapses all terms of a given shape.
\end{question}

\begin{question}[$\scong$ versus $\sim$]\label{q:sim}
The trace compares messages up to $\scong$, which can decohere branches that no
context distinguishes. Coarsening to bisimilarity is available later, at the
price of a fixed-point definition.
\end{question}

\begin{question}[Beyond Clifford]
Completeness of the equations once T-type couplings are present, and the cost
of comparing payloads that carry their own internal sums in R3.
\end{question}

\paragraph{Acknowledgement.} The design was worked out in dialogue with Claude
(Anthropic), which also drafted this note.

\begin{thebibliography}{99}
\bibitem{Amy} M.~Amy. Towards large-scale functional verification of universal
  quantum circuits. \emph{QPL 2018}, EPTCS 287, 2019.
\bibitem{BP} C.~B\u{a}descu and P.~Panangaden. Quantum alternation: prospects
  and problems. \emph{QPL 2015}, EPTCS 195, 2015.
\bibitem{Barandes} J.~A.~Barandes. The stochastic-quantum correspondence.
  arXiv:2302.10778, 2023.
\bibitem{BKMP} D.~E.~Browne, E.~Kashefi, M.~Mhalla and S.~Perdrix. Generalized
  flow and determinism in measurement-based quantum computation. \emph{New J.
  Phys.} 9:250, 2007.
\bibitem{CJPV} T.~Carette, E.~Jeandel, S.~Perdrix and R.~Vilmart. Completeness
  of graphical languages for mixed states quantum mechanics. \emph{ICALP 2019}.
\bibitem{CoeckeDuncan} B.~Coecke and R.~Duncan. Interacting quantum
  observables: categorical algebra and diagrammatics. \emph{New J. Phys.}
  13:043016, 2011.
\bibitem{Dawson} C.~M.~Dawson, H.~L.~Haselgrove, A.~P.~Hines, D.~Mortimer,
  M.~A.~Nielsen and T.~J.~Osborne. Quantum computing and polynomial equations
  over the finite field $\mathbb{Z}_2$. \emph{Quantum Inf. Comput.}
  5(2):102--112, 2005.
\bibitem{DKPW} R.~Duncan, A.~Kissinger, S.~Perdrix and J.~van de Wetering.
  Graph-theoretic simplification of quantum circuits with the ZX-calculus.
  \emph{Quantum} 4:279, 2020.
\bibitem{CQP} S.~J.~Gay and R.~Nagarajan. Communicating quantum processes.
  \emph{POPL 2005}.
\bibitem{JPV} E.~Jeandel, S.~Perdrix and R.~Vilmart. A complete axiomatisation
  of the ZX-calculus for Clifford+T quantum mechanics. \emph{LICS 2018}.
\bibitem{LeiferMilner} J.~J.~Leifer and R.~Milner. Deriving bisimulation
  congruences for reactive systems. \emph{CONCUR 2000}.
\bibitem{MR05} L.~G.~Meredith and M.~Radestock. A reflective higher-order
  calculus. \emph{ENTCS} 141(5):49--67, 2005.
\bibitem{GWv6} L.~G.~Meredith. Races decided by the Born rule (graded-where,
  quantum case, version~6). F1R3FLY-io/publications, \texttt{fuzzyware/}, 2026.
\bibitem{RB} R.~Raussendorf and H.~J.~Briegel. A one-way quantum computer.
  \emph{Phys. Rev. Lett.} 86:5188, 2001.
\bibitem{Selinger} P.~Selinger. Towards a quantum programming language.
  \emph{Math. Struct. Comput. Sci.} 14(4):527--586, 2004.
\bibitem{Vilmart} R.~Vilmart. A near-minimal axiomatisation of ZX-calculus for
  pure qubit quantum mechanics. \emph{LICS 2019}.
\bibitem{YYF} M.~Ying, N.~Yu and Y.~Feng. Alternation in quantum programming:
  from superposition of data to superposition of programs. arXiv:1402.5172,
  2014.
\end{thebibliography}

\end{document}
